%% kapa: exam=1 round=25 year=2014 period=2 form=A subject=appraisal_law subject_ko=부동산관계법규 seq=14 number=14
%% source: Q-Net 감정평가사 기출문제 (cst00309) artlSeq=5101207 file=extracted/1008/5101207/2014년 제25회 감정평가사 1차시험 기출문제/2교시(A형).hwp
%% conversion: kordoc 4.24.51 -> tools/gen1.py (auto); review=proofread-sonnet; flags=none
\begin{question}{14}
국토의 계획 및 이용에 관한 법령상 성장관리방안에 관한 설명으로 옳지 않은 것은?
\begin{choices}
  \item 성장관리방안은 국토교통부장관 또는 시ㆍ도지사가 수립한다.
  \item 도로, 공원 등 기반시설의 배치와 규모에 관한 사항은 성장관리방안에 포함되어야 한다.
  \item 성장관리방안을 수립한 지역에서 하는 개발행위는 중앙도시계획위원회와 지방도시계획위원회의 심의를 거치지 아니한다.
  \item 성장관리방안이 수립된 계획관리지역의 경우 해당 지방자치단체의 조례로 건폐율 및 용적률을 완화하여 적용할 수 있다.
  \item 성장관리방안을 수립하거나 변경한 경우에는 이를 고시하고 일반인이 열람할 수 있도록 하여야 한다.
\end{choices}
\end{question}
\answer{1}
